%%%%%%%%%%%%%%%%%%%%%%%%%%%%%%%%%%%%%%%%%
% Developer CV
% LaTeX Template
% Version 1.0 (28/1/19)
%
% This template originates from:
% http://www.LaTeXTemplates.com
%
% Authors:
% Jan Vorisek (jan@vorisek.me)
% Based on a template by Jan Küster (info@jankuester.com)
% Modified for LaTeX Templates by Vel (vel@LaTeXTemplates.com)
%
% License:
% The MIT License (see included LICENSE file)
%
%%%%%%%%%%%%%%%%%%%%%%%%%%%%%%%%%%%%%%%%%
\documentclass[9pt]{developercv}

% Additional packages for better formatting
\usepackage{enumitem}
\setlist{noitemsep, topsep=0pt, leftmargin=*, nosep}
\usepackage{setspace}
\setstretch{1.05}

\begin{document}

%----------------------------------------------------------------------------------------
% Header - Pixel Perfect Formatting
%----------------------------------------------------------------------------------------

\begin{minipage}[t]{0.49\textwidth}
    \vspace{-\baselineskip}
    \colorbox{black}{{\HUGE\textcolor{white}{\textbf{KRZYSZTOF}}}}\\[0.5ex]
    \colorbox{black}{{\HUGE\textcolor{white}{\textbf{CIEŚLIK}}}}\\[1ex]
    {\huge Embedded Developer}\\
\end{minipage}%
\hfill
\begin{minipage}[t]{0.49\textwidth}
    \vspace{-\baselineskip}
    \icon{MapMarker}{12}{Gdańsk, Pomerania, Poland}\\[1ex]
    \icon{Phone}{12}{REDACTED}\\[1ex]
    \icon{At}{12}{krzysztofcieslik875@gmail.com}\\[1ex]
    \icon{Linkedin}{12}{\href{https://www.linkedin.com/in/krzysztof-cieślik-968128223}{linkedin.com/in/krzysztof-cieślik-968128223}}\\[1ex]
    \icon{Github}{12}{\href{https://github.com/Hathor875}{github.com/Hathor875}}\\[1ex]
\end{minipage}

%----------------------
% Summary - Pixel Perfect Formatting
%----------------------
\cvsect{Summary}

\noindent
I am an electronics and embedded developer with 2 years of professional experience, working mainly
with STM32 bare-metal systems, communication protocols, and PCB design. I combine practical skills in
electronics, automation, and mechatronics with programming in C, C++ and Python, applying them in the
development of new devices, prototypes, and testing tools.

\vspace{0.5ex}
\noindent
I have experience working independently and in R&D teams, including automated hardware validation,
measurement equipment operation, and prototyping (3D printing, small CO laser). My background also
covers hydraulics, pneumatics, and mechanical design from my earlier automation work. I continuously
expand my competencies toward modern embedded software, clean architecture, and reliable engineering
practices.
%----------------------------------------------------------------------------------------
% EXPERIENCE
%----------------------------------------------------------------------------------------

\cvsect{Professional Experience}

\begin{entrylist}
	\entry
		{10.2024--present}
		{Programmer - Electronic Designer}
		{RADMOR S.A., Gdynia}
		{Designing and programming electronic circuits, prototype development, R\&D work}
	\entry
		{07.2024--10.2024}
		{R\&D Department Intern}
		{RADMOR S.A., Gdynia}
		{R\&D team support, electronic equipment testing and documentation}
	\entry
		{02.2024--10.2024}
		{Programming Tutor for children and teenagers}
		{Giganci Programowania, Gdańsk}
		{Teaching programming basics, Python, algorithms}
	\entry
		{12.2023--06.2024}
		{Computer Laboratory Technician}
		{ManpowerGroup, Gdańsk}
		{Computer hardware maintenance and repair, lab administration}
	\entry
		{08.2022--01.2024}
		{Facility Technician}
		{Sodexo, Gdańsk}
		{Electrical installation maintenance and repair, technical facility support in a skyscraper}
	\entry
		{07.2022--09.2022}
		{Automation Technician}
		{Ceramika Paradyż, Opoczno}
		{PLC controller programming, IT support and maintenance}
	\entry
		{06.2021--09.2021}
		{Automation Technician}
		{CERSANIT GROUP, Opoczno}
		{Production machine maintenance, process continuity assurance}
	\entry
		{07.2020--09.2020}
		{Electrician}
		{PKP Intercity Remtrak Sp. z o.o., Opoczno}
		{Repair of electrical systems in passenger railway cars}
\end{entrylist}

% --- END OF EDUCATION SECTION ---

%----------------------------------------------------------------------------------------
% ADDITIONAL INFORMATION
%----------------------------------------------------------------------------------------

\cvsect{Key Skills}

% Two columns of skills with bullet points
% Left column
\begin{minipage}[t]{0.48\textwidth}
\begin{itemize}
  \item Python
  \item Electronics
  \item Electrical Engineering
  \item Automation
  \item Embedded Systems
  \item Communication Protocols (i2c,spi,eth)
\end{itemize}
\end{minipage}%
\hfill
% Right column
\begin{minipage}[t]{0.48\textwidth}
\begin{itemize}
  \item Automated Validation Hardware (Python)
  \item PCB Design and circuit
  \item Microcontroller Programming (mainly stm32)
  \item Equipment Diagnostics and Repair
  \item C/C++
  \item Linux
\end{itemize}
\end{minipage}

\vspace{0.5cm}

\begin{minipage}[t]{0.3\textwidth}
	\vspace{-\baselineskip}
	\cvsect{Languages}
	
	\begin{itemize}
		\item Polish – native
		\item English – B1
	\end{itemize}
\end{minipage}

%----------------------------------------------------------------------------------------
% DATA PROCESSING CONSENT FOOTER
%----------------------------------------------------------------------------------------

\vfill


\clearpage

\cvsect{More About Me}

I am an embedded developer with a multidisciplinary background that combines electronics,
automation, software engineering, and mechanical design. I enjoy working on projects where
hardware and software intersect, and where thoughtful engineering decisions make systems
not only functional, but reliable, understandable, and maintainable.

My earlier experience in industrial automation shaped my engineering mindset: precision,
robustness, and problem-solving under pressure. This foundation, combined with embedded
software development and prototyping skills, influences the way I design systems today.
I value clarity, predictable behavior, and solutions that scale over time rather than
quick fixes that create technical debt.

Beyond electronics and programming, I work with 3D printing, mechanical modelling
(FreeCAD, OpenSCAD), and small CO laser machining. I also build physical prototypes
and test rigs, solder SMD/BGA components, and operate lab equipment such as oscilloscopes,
signal generators, power analyzers, and VNA instruments. This hands-on experience helps
me design systems that account for physical reality, not only theory.

What motivates me most is creating tools, devices, and workflows that make engineering
smoother, clearer, and more enjoyable for the people who use them—whether it’s a team,
future maintainers, or my future self.


\cvsect{Engineering Philosophy}

I believe good engineering is a craft. It requires intention, documentation, and thoughtful
architecture. My approach focuses on designing systems that are easy to extend, test, and
reason about. Instead of relying on shortcuts or tightly coupled designs, I invest in
structure and clarity—from consistent workflows and CI/CD to well-documented interfaces
and reproducible development environments.

Clean architecture, maintainability, traceability, and solid engineering practices
(SOLID, modular design, layered firmware structure) guide my work. I strive to build
solutions that “ratchet forward”: each iteration should make the system stronger, cleaner,
and easier to maintain.


\cvsect{Engineering Thesis Project}

As the leader of a five-person engineering team, I am currently developing a system for
creating structured technical documents with templating, metadata management, and automatic
versioning. The project is built in Python (FastAPI) with a React-based GUI running inside
Tauri, allowing cross-platform distribution and a native-application feel.

My responsibilities include:

– architecture design and defining the entire toolchain  
– CI/CD pipeline development  
– workflow automation and project structure  
– documentation standards and requirement gathering  
– internal team tools (time tracking, communication utilities)  
– maintaining the local hosting environment for services  

This project reflects how I work: treating software as a well-designed system rather than
a collection of scripts. I place strong emphasis on clear separation of concerns,
maintainability, and reproducibility so the application can evolve with the team’s needs.


\cvsect{Why This Page Exists}

A résumé lists skills and history.  
This page explains how I think, how I approach engineering problems, and what drives
my work. If you want a deeper understanding of what I bring to a project or a team,
this narrative is designed to provide that context.

\end{document}
